\documentclass[10pt,letterpaper]{article}
\usepackage{cornell-notes}

\title{Clause 05.02.01: Context Of Architecture Description}
\author{}
\date{}
\setDocTitle{Clause 05.02.01: Context Of Architecture Description}
\setDocSubtitle{ISO/IEC/IEEE 42010:2022}
\setDocOwner{ISO 42010 Cornell Notes}
\setDocDate{\today}
\setCornellCollection{ISO/IEC/IEEE 42010:2022 Cornell Notes}
\setCornellUnitType{Clause}
\setCornellUnitNumber{05.02.01}
\setCornellUnitTitle{Context Of Architecture Description}
\hypersetup{pdftitle={Clause 05.02.01: Context Of Architecture Description},pdfsubject={Cornell notes},pdfauthor={}}

\begin{document}
\maketitle
\begin{CornellOverviewBox}{Overview}
{\Large\bfseries\textcolor{CNavy}{Section 5.2.1: Context of architecture description}}\\[3pt]
{\small\textbf{Classification:} Conceptual foundation}\\
{\small\textbf{Learning objective:} Explain the role of context of architecture description and connect it to related architecture-description concepts.}
\end{CornellOverviewBox}
\vspace{3pt}

\begin{CornellNotesTable}
\CornellNoteRow{Purpose / key concept}{The term "entity of interest" is used Here to refer to the subject of an architecture description.}
\CornellNoteRow{Core details}{\begin{itemize}[leftmargin=*,nosep,topsep=1pt]
\item The term is intended to encompass, but is not limited to, entities within the following fields of application, reflecting the intended scope of This section as specified in Clause 1.
\end{itemize}}
\CornellNoteRow{Example / application}{The identification of an EoI generally results from the definition of the problem space. The problem description can be expressed with an architecture definition of a set of operational capabilities (often called “capability architecture”). At this capability definition stage, identified stakeholders are potentially concerned by the future EoI.}
\CornellNoteRow{Interpretive note}{\begin{itemize}[leftmargin=*,nosep,topsep=1pt]
\item The figures and text in the remainder of Clause 5 constitute a set of conceptual models of architecture description. Figures 1 to 6 use an informal entity-relationship diagram notation to facilitate comprehension by users of This section. In the figures, rounded rectangles represent information objects, and arrows represent relationships between objects with the annotation read in the arrow direction. The figures illustrate the key concepts described throughout Clause 5. Annex A presents the full conceptual model.
\item Identification of the EoI can emerge from the analysis of the concerns of the stakeholders or can preexist before identification of some stakeholders and their concerns.
\end{itemize}}
\CornellNoteRow{Active recall}{\begin{itemize}[leftmargin=*,nosep,topsep=1pt]
\item What is the central role of Context of architecture description?
\item How does Context of architecture description connect to adjacent architecture-description concepts?
\end{itemize}}
\end{CornellNotesTable}


\begin{CornellSummaryBox}{Summary}
The term "entity of interest" is used Here to refer to the subject of an architecture description.
\end{CornellSummaryBox}

\end{document}
